\documentclass[11pt]{article}
\usepackage[margin=1in]{geometry}
\usepackage{hyperref}
\usepackage{url}
\title{TLA+ Study Group — Week 03}
\author{TLA+ Study Group}
\date{\today}
\begin{document}
\maketitle

\section*{Summary}
% bullets
\begin{description}
    \item [Introduction to TLA+] TLA+ marketing and reasoning
    \item [State Machines in TLA+] First intro to a specific state machine and glimpse of TLA+
    \item [Resources and Tools] TLA+ Toolbox walkthrough
    \item [Die Hard] 3 liter jug + 5 liter jug, get 4 liters problem walkthrough
\end{description}

\section*{Good stuff}
% bullets
\begin{itemize}
    \item TLA in principle reminds me of kripke models used in temporal and dynamic epistemic logic\footnote{\url{https://plato.stanford.edu/entries/dynamic-epistemic/appendix-A-kripke.html}},
    but so far Lamport just uses the full Kripke models as illustrative, with no mention of those.
    \item Equivalence vs assignment issue is real. Getting used to treating $=$ as just a part of formula, with assignment dictated by prime ($'$)
    versions of the variables is counterintuitive.
    \item MapReduce is essentially encoded in the logic behind array assignments (made even funnier by \LaTeX's \texttt{mapsto}): $$[i \in 1..42 \mapsto i^2]$$
\end{itemize}

\section*{Weird stuff}
\begin{itemize}
    \item The equvalence vs assignemnt issue also means that you have to consider the state of ALL variables in every state described. If
    operation needs to be performed on just one variable - the rest still need to be included in the state description; at least the ones
    that have to remain invariant. If this is not the case, formulas suffer.
    \item TLA+ Toolbox doesn't really vibe well with Sway + Wayland setup. Crashed on me thrice.
    \item Dumping conventions and using \texttt{/=} instead of \texttt{!=} for $\neq$ is cruel, but drives the point of ``this is math, not programming''
    home pretty well.
\end{itemize}


% bullets

\section*{Open Questions}
% bullets

\end{document}
